\section{总结}\label{sec:conclusion}
本文围绕GraphMAE的节点特征掩码、GAT编码、重掩码解码与SCE目标，完成了Cora上的课程级部分复现。官方PyG固定版本在五种子下得到83.98\%的测试均值和1.056882个百分点的样本标准差，与原论文84.2\%的报告均值接近，但初始化次数与环境差异使其不能视为完全相同协议的复现。

在自主实验中，本文只改变掩码节点采样，设计并实现固定预算的度数加权无放回策略。二十次训练完成后，验证集在三个正向系数中选择0.25；最终其测试均值82.92\%低于同协议均匀控制组84.20\%，平均配对变化为$-1.28$个百分点。实际掩码统计证明高度数节点被更频繁选中，但三类度数组的平均分类准确率均下降，末100轮SCE也更高。

这些结果表明，增加某类节点的重建机会不必然改善下游表示。本文的贡献在于对一个明确、可实现的采样假设进行受控检验并保留负结果，而非提出已证实有效的新算法。结论受单数据集、少量种子、有限参数候选及数值非确定性限制，相关机制解释仍需独立证据支持。
